\documentclass[conference]{IEEEtran}
\usepackage[utf8]{inputenc}
\usepackage[T1]{fontenc}
\usepackage{cite}
\usepackage{amsmath,amssymb}
\usepackage{graphicx}
\usepackage{booktabs}
\usepackage{url}

\title{Latency‑Aware Selective Verification in Generate--Validate NetOps Pipelines: A Preregistered Paired Study with gpt‑5‑mini}

\author{
\IEEEauthorblockN{Autonomous AI Researcher}
\IEEEauthorblockA{CNSM 2026 Agentic AI Researcher Track}
}

\begin{document}

\maketitle

\begin{abstract}
We preregistered (CAND‑2026‑001) and executed a paired, artifact‑backed experiment that measures whether inserting a selective low‑latency verification tier (risk scoring + fast validator) into a generate--validate--repair NetOps pipeline reduces unsafe commits while satisfying a preregistered latency SLO. Using a deterministic synthetic intent corpus (N=400 paired intents; 10\% deterministic adversarial variants), gpt‑5‑mini as the generator, and a deterministic full checker (deterministic\_netops\_validator\_v5) as ground truth, we observed baseline valid proposals = 399/400 (0.9975) and guarded valid proposals = 400/400 (1.000). Paired difference (guarded - baseline) = 0.0025; exact McNemar two‑sided p = 1.0; paired‑bootstrap 95\% percentile CI = [0.0, 0.0075] (10,000 resamples, seed=1729). The preregistered latency hypothesis (H2) could not be evaluated because the required per‑stage timing telemetry (generation, risk scoring, fast validator, full checker timestamps) was not archived in the supplied artifact bundle (see Disclosure Statement and Execution Manifest). All principal artifacts (raw responses, per‑episode scorer JSONs, validator traces, provider‑call logs, model configuration, analysis artifacts and hashes) are archived to permit deterministic replay and independent inspection.
\end{abstract}

\section{Introduction}
Automating intent\ensuremath{\rightarrow}configuration translation is central to modern NetOps but raises two operational constraints: safety (avoid outages/regressions) and latency (meet human and system SLOs). Recent work integrates LLMs into generate--validate--repair and agentic NetOps pipelines, often pairing generators with heavyweight validators; this improves correctness but may violate latency budgets. We preregistered a confirmatory paired experiment (CAND‑2026‑001) to test whether a selective, low‑latency verification tier (risk scoring + fast validator) reduces unsafe commits relative to LLM‑only while meeting a 95th‑percentile latency SLO. The run is fully artifacted: generator configuration, per‑episode scorer JSONs, validator traces, provider‑call logs, execution and analysis manifests, and cryptographic hashes are archived to enable deterministic reproduction of the primary numerical claims.

\section{Related Work}
Prior work documents LLMs for intent translation, configuration generation, and configuration repair, and separately shows that verifier‑augmented pipelines reduce regressions (e.g., generate--validate--repair architectures and benchmarks that combine LLMs with formal/network verification). Unlike benchmark style leaderboards, our contribution is a preregistered paired design that isolates the verifier's effect conditional on the same generated candidate set (shared‑initial protocol). We position our study as an operationally traceable measurement of the marginal verifier effect under a tightly controlled, reproducible workflow; nearest‑work artifacts used for study design are archived with this run (see Cited Record Ids).

\section{Methodology}
Preregistration and design. The study preregister (CAND‑2026‑001) froze: research question; primary estimand = paired\_success\_rate\_difference\_guarded\_minus\_baseline; confirmatory hypotheses H1 (safety) and H2 (latency); N=400 paired intents with deterministic 10\% adversarial variants; shared\_initial protocol (independent\_condition\_generation=false); retry/missingness rules; analysis executor paired\_binary\_analysis\_v1; and bootstrap parameters (10,000 resamples, seed=1729). The preregistration manifest is recorded in the execution manifest (preregistration\_sha256: 086d248e89636e99e788232e78a99b63459648545e6e79ed763c29521f8f8572).
Corpus and tools. Confirmatory corpus: deterministic synthetic intent generator (400 unique intents; 40 deterministic adversarial items). Generator: gpt‑5‑mini (execution/model\_configuration.json, SHA‑256: a40e96e212ab87da251ac0d6dbb75bc543afa13438b5a6b9ebd073597ffdd820). Validator / full checker: deterministic\_netops\_validator\_v5; per‑episode scorer JSONs are archived at execution/scoring/*.json. Adapter/orchestration family: hosted\_netops\_gvr\_v1; provider call traces are archived under execution/provider\_calls/*.json (see Execution Manifest).
Shared‑initial protocol and estimand. We preregistered shared\_initial candidates for both conditions to isolate verifier/repair contributions given the same generator output (independent\_condition\_generation=false). This produces a conditional estimand: the verifier effect given the observed generated candidate set. The reconciliation artifact documents model\_calls\_used=401 (shared\_initial\_generation=400 + repair=1) versus the planned 800 model calls under independent generation (analysis/deterministic\_reconciliation.json; SHA‑256: a7570bf6ff3a946e...).
Scoring and statistical tests. Each candidate was deterministic‑scored by the archived validator (binary valid/invalid plus structured trace). Primary test: exact McNemar two‑sided on the paired contingency (n11,n10,n01,n00) as implemented in paired\_binary\_analysis\_v1 (archived). Confidence interval: paired bootstrap (percentile) with 10,000 resamples and seed=1729. All scoring outputs and the contingency table are archived (analysis/paired\_contingency\_table.csv SHA‑256: a6906bb43712dcf8...).

\section{Results}
Execution and accounting. The run completed preregistered confirmatory set: completed\_episode\_count = 800 (400 pairs), model\_calls\_used = 401 (see analysis/deterministic\_reconciliation.json SHA‑256: a7570bf6ff3...), and all raw responses, provider calls, and per‑episode scorer JSONs are archived (execution/*; hashes in the Execution Manifest).
Confirmatory safety (H1). The archived paired contingency (analysis/paired\_contingency\_table.csv; SHA‑256: a6906bb4...) contains n11=399, n10=1, n01=0, n00=0. Thus baseline valid proposals = 399/400 = 0.9975; guarded valid proposals = 400/400 = 1.000. Paired difference (guarded - baseline) = 0.0025. The exact McNemar two‑sided test implemented in the archived analysis executor returned p = 1.0 for discordant counts (n10=1,n01=0); the executor and invocation parameters are archived (analysis/analysis\_log.jsonl; SHA‑256: 9ea0d56d...). Paired‑bootstrap (percentile) 95\% CI = [0.0, 0.0075] (10,000 resamples, seed=1729).
Why p = 1.0 in the archived exact two‑sided McNemar implementation. The analysis used the exact two‑sided definition archived in paired\_binary\_analysis\_v1; for extremely small discordant counts the implementation's two‑sided tail mass assignment yields p=1.0 in the observed asymmetric case (n10=1,n01=0). The exact implementation and parameters (bootstrap seed, resamples) are archived so readers can re‑run the identical decision rule and obtain the same p‑value (artifact paths and hashes listed below).
Representative inspectable evidence. The single discordant pair is archived for deterministic inspection at these exact paths and SHA‑256s: execution/scoring/task-000149-baseline.json (SHA‑256: fe147f43e33432ed446d76500c2c9566ca5f72086abbacb5ee84b13774a9d18e) and execution/scoring/task-000149-guarded.json (SHA‑256: d0cbfe9a4181cd87f35c9ead5b6697e1adcfde23e52461e8c21e5219fe607ca3). Each per‑episode JSON contains: normalized candidate, raw candidate hash, validator trace (parsed commands, any violations, final state), and provider‑call references; these files enable readers to replay why the baseline candidate was scored invalid while the guarded pipeline accepted the same (or repaired) candidate.
Latency (H2) status. The study preregistered a latency hypothesis (H2): 95th‑percentile end‑to‑end pipeline latency <= 2000 ms for the selective verifier and a preregistered paired comparison of 95th‑percentiles. The artifact bundle does not contain the per‑stage timing telemetry required to evaluate H2. In particular, per‑stage timestamps for generation, risk scoring, fast validator, and full checker were not archived in the execution/provider\_calls/*.json nor in execution/execution\_log.jsonl; the absence is documented in the execution manifest. Therefore H2 is unresolved for this run and no latency claims are made.

\section{Discussion}
Statistical and operational interpretation. The observed absolute safety gain is numerically small because the generator produced near‑ceiling valid candidates on our synthetic corpus (baseline 0.9975). The paired design with shared\_initial candidates isolates verifier contribution given that candidate set; it does not measure the joint generator+verifier effect that would arise if each condition used independent generator sampling. For deployments where generator baseline is lower or where adversarial items are common, the verifier could yield larger absolute gains. The preregistration included exploratory scans (verification thresholds and adversarial subgroup analyses) and the full per‑episode scored artifact set permits those analyses.
Design tradeoffs. The shared\_initial (conditional) estimand reduces variance due to generator sampling and yields fewer model calls (401 used versus 800 under independent generation). This choice makes the result interpretable as the marginal verifier effect for the observed candidates; practitioners must choose the estimand (conditional vs joint) aligned with operational decisions. Analysis/deterministic\_reconciliation.json documents the exact accounting (model\_calls\_used=401; SHA‑256: a7570bf6ff3a...).
Limitations and threats to validity. Construct validity: the deterministic\_netops\_validator\_v5 is a modeling choice with finite semantic fidelity to vendor devices and digital twins. External validity: single model (gpt‑5‑mini), synthetic corpus, and single deterministic validator limit generalization. Internal validity: the shared\_initial design was preregistered and reconciled; no undisclosed reruns are present per the execution manifest. Conclusion validity: the small discordant counts produce wide uncertainty and an exact McNemar p that depends on the archived exact two‑sided definition; readers can re‑execute the same test using archived analysis artifacts.
Operational recommendation. For operational adoption: (1) measure generator baseline on the real intent distribution; (2) evaluate the selective verifier's false‑negative rate on OOD/adversarial intents; (3) ensure per‑stage latency telemetry is recorded and archived so latency SLOs (like H2) can be evaluated as preregistered. Our artifact bundle documents how to replay the run and inspect any episode deterministically.

\section{Limitations}
\begin{itemize}
\item Synthetic deterministic corpus and single generator (gpt‑5‑mini): external validity to heterogeneous real operator intent distributions and other LLMs is limited.
\item Single deterministic validator (deterministic\_netops\_validator\_v5): construct validity depends on validator semantic fidelity to vendor devices and production digital twins.
\item Shared‑initial design (independent\_condition\_generation=false) yields a conditional estimand (verifier effect given the observed generated candidates) rather than the joint generator+verifier estimand; model\_calls\_used=401 vs planned\_model\_calls=800 is reconciled in analysis/deterministic\_reconciliation.json (SHA‑256: a7570bf6...).
\item Latency hypothesis (H2) could not be evaluated because required per‑stage timing telemetry (generation, risk scoring, fast validator, full checker timestamps) was not archived; no preregistered latency claim is made.
\end{itemize}

\section{Conclusion}
We ran a fully preregistered, artifacted paired experiment to measure the marginal safety impact of a selective low‑latency verification tier in an LLM\ensuremath{\rightarrow}NetOps pipeline. The guarded verifier produced one additional accepted valid proposal in N=400 paired intents (paired difference = 0.0025; exact McNemar two‑sided p=1.0; paired bootstrap 95\% CI = [0.0,0.0075], 10,000 resamples, seed=1729). The latency hypothesis (H2) cannot be evaluated because required per‑stage timing telemetry was not archived; this absence is documented and must be remedied in future confirmatory runs. All principal artifacts (raw responses, per‑episode scoring JSONs, validator traces, provider calls, model configuration, execution manifest, and analysis artifacts) are archived with SHA‑256 hashes and file paths to enable deterministic reproduction and independent inspection of the single discordant example described above. Future work should evaluate multiple models, validators, and independent‑generation designs, and must archive per‑stage latency telemetry for operational latency claims.

\section*{Disclosure Statement}
Disclosure Statement (mandatory, artifact‑backed):
- Initial master prompt (immutable archived reference): provenance/master\_prompt.txt (SHA‑256: 1872df1e1805d2d96940456ca016bd665d1d5196add77f5acdf1582bb39b15ba). The Research Director selected the topic family (Generative AI / LLMs for NetOps) and approved pre‑lock constraints. After the autonomy lock the autonomous principal researcher performed literature review, study selection, preregistration, experiment execution, analysis, manuscript drafting, internal AI peer review, and revisions without manual post‑lock editing of technical content.
- Preregistration: Study ID CAND‑2026‑001 (frozen prior to execution). The preregistered plan (recorded in the execution manifest) specified the primary estimand (paired\_success\_rate\_difference\_guarded\_minus\_baseline), confirmatory hypotheses H1 and H2, N=400 paired intents with 10\% deterministic adversarial variants, shared\_initial protocol (independent\_condition\_generation=false), retry/missingness rules, analysis executor (paired\_binary\_analysis\_v1), and bootstrap parameters (10,000 resamples, seed=1729). Preregistration SHA‑256: 086d248e89636e99e788232e78a99b63459648545e6e79ed763c29521f8f8572 (execution manifest).
- Models, adapter, orchestration: Generator declared = gpt‑5‑mini; model configuration archived at execution/model\_configuration.json (SHA‑256: a40e96e212ab87da251ac0d6dbb75bc543afa13438b5a6b9ebd073597ffdd820). Adapter/orchestration family: hosted\_netops\_gvr\_v1. Provider call traces and per‑call artifacts are archived under execution/provider\_calls/*.json (full hashes in the execution manifest). The run used 401 provider model calls (shared\_initial\_generation=400 + repair=1); planned\_model\_calls per preregistration = 800. The reconciliation artifact documents model\_calls\_used=401 (analysis/deterministic\_reconciliation.json; SHA‑256: a7570bf6ff3a946e8397f02452ab846c82d5184b7768f006bd06359a1b177de7).
- Validators and scoring: Ground truth / full checker = deterministic\_netops\_validator\_v5; archived per‑episode scorer JSONs are available at execution/scoring/*.json. Principal analysis artifacts: analysis/paired\_contingency\_table.csv (SHA‑256: a6906bb43712dcf82b3bc8bdc0ef095d228bcdc8a6a857d9cbbba1ad87245af8), analysis/analysis\_log.jsonl (SHA‑256: 9ea0d56d48a5cc7ff198d85d0a12c21536fab716c042bd5eb2033492eda5b18e), analysis/deterministic\_reconciliation.json (SHA‑256: a7570bf6ff3a946e8397f02452ab846c82d5184b7768f006bd06359a1b177de7). The single discordant episode that underpins the primary numerical claim is archived at: execution/scoring/task-000149-baseline.json (SHA‑256: fe147f43e33432ed446d76500c2c9566ca5f72086abbacb5ee84b13774a9d18e) and execution/scoring/task-000149-guarded.json (SHA‑256: d0cbfe9a4181cd87f35c9ead5b6697e1adcfde23e52461e8c21e5219fe607ca3). Each per‑episode JSON contains the normalized candidate, raw candidate hash, validator trace (parsed commands, violations, final state), and provider‑call references to enable deterministic replay.
- Statistical implementation: The confirmatory analysis used the archived paired\_binary\_analysis\_v1 executor. Primary test: exact McNemar (two‑sided) on the paired contingency (n11=399,n10=1,n01=0,n00=0); the archived exact two‑sided implementation returned p=1.0 for the observed asymmetric discordant counts. Confidence interval computed via paired bootstrap (percentile), 10,000 resamples, seed=1729 (parameters archived in analysis/analysis\_log.jsonl).
- Missing archived latency telemetry: The preregistration required per‑stage latency telemetry (generation, risk scoring, fast validator, full checker timestamps) to evaluate H2. These per‑stage timing fields were not archived in the supplied artifact bundle (they are absent from execution/provider\_calls/*.json and execution/execution\_log.jsonl); the absence is documented in the execution manifest and analysis/deterministic\_reconciliation.json. Consequently the preregistered latency hypothesis (H2) is unresolved for this run and no latency claims are made.
- Human involvement and autonomy boundary: Humans performed pre‑lock infrastructure development, topic selection, and preregistration constraint approval. After the autonomy lock the autonomous system conducted literature review, study selection, preregistration, experiment execution, analysis, manuscript drafting, internal AI peer review, and finalization without manual post‑lock editing of technical content. All retries, failures, and deviations are logged in the execution manifest.
- Reproducibility: All principal outputs (raw responses, per‑episode scorer JSONs, execution/provider\_calls, model configuration, execution manifest, analysis artifacts and logs) are archived with SHA‑256 hashes (see execution\_manifest) so third parties can deterministically reproduce and inspect the run. The exact artifact file paths and SHA‑256s referenced above enable direct deterministic inspection of the key discordant episode and the primary analysis artifacts.

\begin{thebibliography}{99}
\bibitem{ref1} Zhenrong Gu, Peng Zhang, Xing Feng, Xu Liu, "Astragalus: Automatic Configuration Repair for Production Networks", 2026
\bibitem{ref2} Jordan Augé, Sam Betts, Giovanna Carofiglio, Giulio Grassi, Martin Gysi, John Kenneth d'Souza, "Aether: Network Validation Using Agentic AI and Digital Twin", 2026
\bibitem{ref3} Ioannis Protogeros, Rufat Asadli, Benjamin Hoffman, Laurent Vanbever, "Benchmarking LLM-Driven Network Configuration Repair", 2026, 10.3929/ethz-c-000799709
\bibitem{ref4} Abdelkader Mekrache, Adlen Ksentini, "LLM-enabled Intent-driven Service Configuration for Next Generation Networks", 2024, 10.1109/netsoft60951.2024.10588881
\bibitem{ref5} Jiayi Mao, Liqun Li, Yanjie Gao, Zegang Peng, Shilin He, Chaoyun Zhang, Si Qin, S Khairou Khalid, Qingwei Lin, Saravan Rajmohan, Sitaram Lanka, Dongmei Zhang, "StepFly: Agentic Troubleshooting Guide Automation for Incident Diagnosis", 2026, 10.1145/3808143
\bibitem{ref6} Changhua Pei, Zexin Wang, Fengrui Liu, Zeyan Li, Yang Liu, Xiao He, Rong Kang, Tieying Zhang, Jianjun Chen, Jianhui Li, Gaogang Xie, Dan Pei, "Flow-of-Action: SOP Enhanced LLM-Based Multi-Agent System for Root Cause Analysis", 2025, 10.1145/3701716.3715225
\bibitem{ref7} https://doi.org/10.1109/tnet.2023.3269983
\end{thebibliography}

\end{document}
